While Large Language Models (LLMs) have demonstrated impressive capabilities in spatial reasoning and pathfinding benchmarks, the linguistic structure of their generative spatial narratives remains under-explored. We investigate whether LLMs exhibit a distinct "spatial grammar" when describing movement through physical environments. By prompting \gpto and \gptom to generate narratives based on layouts from the \babitask dataset, we analyze linguistic features including spatial preposition entropy, sentence length variance, and Part-of-Speech (POS) bigram distributions. Our findings reveal that default LLM generations are highly formulaic, characterized by low spatial vocabulary entropy (3.41 bits for \gpto) and rigid, repetitive sentence structures. However, we show that explicit prompting for human-like variance significantly increases structural diversity and introduces descriptive sensory details. These results suggest that the "RLHF tone" induces a systematic bias in how LLMs represent physical space, favoring topological mapping over naturalistic storytelling.
